\documentclass[12pt,letterpaper]{article}
\usepackage{amsmath,amssymb}
\listfiles
\paperwidth=8.5in
\paperheight=11in
\textwidth=6.5in
\textheight=9in
\oddsidemargin=0pt
\evensidemargin=0pt
\topmargin=0pt
\headheight=0pt
\headsep=0pt
\footskip=0pt
\parindent=0pt
\parskip=12pt
\pagestyle{empty}
\begin{document}
\begin{minipage}{240pt}
A student writes a careful explanation of an equation. The paragraph must justify ordinary spaces and choose reproducible line breaks. Repeated words expose changes in font widths: minimum maximum minimum maximum. Hyphenation illustrates international collaboration and computational mathematics.

A second paragraph starts here. Its baseline and separation must remain distinct. Short final lines do not stretch like fully justified lines.
\end{minipage}

\begin{minipage}{180pt}
The same font is constrained to a narrower measure. Natural line breaking must change without moving the paragraph's left edge.
\end{minipage}
\end{document}
